\chapter{Mode 3: Orthogonal Backgrounds over an Isometric One}
\label{ch:mode3}

\section{Overview}

Mode 3 keeps one isometric background and gives the fetch slots of the other to two orthogonal ones. It is for a game
with an isometric ground and a user interface that needs more than the status layer of Mode 2: a status display and a
layer of windows over it, with the map, and the vehicles and buildings as sprites, underneath.

\begin{description}
	\item[Selected by:] \texttt{MODE} = 2 in the Base Configuration register.
	\item[Layers:] Orthogonal backgrounds 1 and 2 in front, and an isometric background 3 behind them. The order front to back is
	      S1, B1, S2, B2, S3, B3, as in Mode 1 (Section~\ref{sec:m1-priority}).
	\item[Registers:] Each background uses its own tilemap, tile data, scroll and palette configuration registers, as in Mode 1.
	      Background 3 is the isometric one, so \texttt{B3MAP}, \texttt{B3TILES}, \texttt{B3SCRX}, \texttt{B3SCRY} and
	      \texttt{B3PAL} follow Mode 2 (Chapter~\ref{ch:mode2}).
\end{description}

Mode 3 adds no hardware to Mode 2. Each background is the same fetch unit as in the other modes, and the mode only says
which slots of a group of four pixels belong to which background. The price is the status layer of Mode 2: with no
spare slots there is no queue, and no second isometric background.

\section{The Orthogonal Backgrounds}

An orthogonal background is the background of Mode 1 (Chapter~\ref{ch:mode1}): a tilemap of 8 by 8 or 16 by 16 tiles,
32 by 32 entries or larger (Section~\ref{sec:tilemap-size}), that wraps in both directions. Its tilemap entry holds a
tile number and a palette selection (Section~\ref{sec:tilemap-entry}), and its flip bits and priority bit apply as they
do in Mode 1. Tiles are 4 bits per pixel, 16 words each (Section~\ref{sec:m1-tile-layout}), or, when \texttt{T2} is set
in its palette configuration register, 2 bits per pixel, 8 words each (Section~\ref{sec:flat-two-bit}). The palette is
chosen as in Section~\ref{sec:palette-select}, with the register's own \texttt{OFFSET}. Pixels of value 0 are
transparent. The priority bit puts a tile in front of every sprite layer.

\subsection{Scrolling}
\label{sec:m3-scrolling}

Each orthogonal background scrolls independently. Its scroll registers are sampled once per line, as in Mode 2. The
device fetches groups of four pixels aligned to the plane, so that group $g$ of the line is the plane group $\lfloor
	\texttt{SCRX}/4 \rfloor + g$. Plane group $q$ is the left half of tile column $\lfloor q/2 \rfloor \bmod 32$ when $q$
is even and the right half when it is odd, and the row is $(\text{line} + \texttt{SCRY}) \bmod 256$. When \texttt{SCRX}
is not a multiple of four, the line's first pixel is at position $\texttt{SCRX} \bmod 4$ in a window of two groups, as
in Section~\ref{sec:m2-fetch}. Fine scrolling costs no extra fetches.

\section{Fetch Schedule}

The eight slots of a group of four pixels have fixed owners:

\begin{table}[H]
	\centering
	\begin{tabular}{cp{5.8cm}p{5.2cm}}
		\toprule
		\textbf{Slot} & \textbf{Fetch}                                                & \textbf{Used}                            \\
		\midrule
		0             & Background 1 tilemap entry                                    & Always (2-bit tiles: every second group) \\
		1             & Background 1 tile data word                                   & Always (2-bit tiles: every second group) \\
		2, 3          & The same two fetches for background 2                         &                                          \\
		4--7          & The four fetches of an isometric background, for background 3 & As in Table~\ref{tab:m2-slots}           \\
		\bottomrule
	\end{tabular}
	\caption{Mode 3 Fetch Slots per Group of Four Pixels}
	\label{tab:m3-slots}
\end{table}

The orthogonal backgrounds use four slots in every group, and the isometric background uses two to four, so a group
uses six or more of the eight and the worst case uses all eight. No slot is left for anything else in the visible part
of a line. Sprites are fetched in the horizontal blanking interval, as in the other modes (Chapter~\ref{ch:sprites}),
and the back-porch prefetch of two groups is the same as in Mode 2.

Because no background depends on a slot that another leaves spare, the schedule is fixed and needs no queue.

\section{Differences from Mode 2}

\begin{itemize}
	\item There is no status layer: the orthogonal backgrounds take its place, and \texttt{B3D} disables the isometric
	      background.
	\item \texttt{B1D}, \texttt{B2D} and \texttt{B3D} disable the corresponding background, as in Mode 2.
	\item The isometric background is background 3, so it is behind the orthogonal ones. The shading bits of its tilemap entries
	      (Section~\ref{sec:m2-shading}) apply to it; the orthogonal backgrounds ignore the flip bits.
	\item Only one background is isometric, so the height offsets of Section~\ref{sec:m2-layers} between two isometric
	      backgrounds are not available. Buildings and vehicles are sprites.
\end{itemize}

\section{Restrictions}

\begin{itemize}
	\item One isometric background, with the restrictions of Mode 2 (Section~\ref{sec:m2-restrictions}).
\end{itemize}

\section{Programming Mode 3}

\begin{enumerate}
	\item Write the tilemap, tile data and palettes of the three backgrounds. The orthogonal tiles use the format of
	      Section~\ref{sec:m1-tile-layout} or~\ref{sec:flat-two-bit}; the isometric tiles use the staircase format
	      (Section~\ref{sec:m2-tile-format}) or its 2-bit form.
	\item Set \texttt{B1MAP}--\texttt{B3MAP}, \texttt{B1TILES}--\texttt{B3TILES} and the palette and scroll registers.
	\item Write 2 to the \texttt{MODE} field. Do it in vertical blanking: a change during a line takes effect at the start of the
	      next one.
\end{enumerate}
